\section{Conclusion}
\label{sec:conclusion}

CT-KAT binds three checkers --- a Valgrind/Memcheck structural CT
check~\cite{valgrind,ctgrind}, a warn-only variable-latency asm-scan, and the
dudect statistical timing test~\cite{dudect} --- under one YAML config and
compiler-by-cflags matrix, behind a default-deny triage taxonomy with a cited
accepted-variable-time registry. On the corpus it is
\emph{layer-justified} (each layer single-covers a leak the others miss,
Table~\ref{tab:coverage}), \emph{validated} (verdict classes map to concrete
targets, Tables~\ref{tab:corpus} and~\ref{tab:dudect}), and \emph{honest by
construction}: on ML-DSA-65~\cite{fips204} the taxonomy refused to auto-label a
structural failure as key-recovery, holding the unregistered leak sites at
\emph{needs-analysis} rather than over-claiming. The KyberSlash
case~\cite{kyberslash} shows structural checking alone is insufficient for PQC,
justifying the multi-layer design.

We claim no new attack, guarantee, or absence proof; the phenomena are known,
and the contribution is their integration and triage. Future work broadens the
corpus to Falcon and SPHINCS+ and resolves ML-DSA attribution via a dedicated
\texttt{-fno-inline} build.
